% !PRINTED 442
\weberpage{442}
%<<CONT

\begin{tabular}{@{}l l l l@{}}
1. & Substitutions with three poles, & & \\
2. & \textquotedbl & \textquotedbl & two poles,\\
3. & \textquotedbl & \textquotedbl & one pole,\\
4. & \textquotedbl & \textquotedbl & one principal axis and one pole,\\
5. & \textquotedbl & \textquotedbl & one principal axis.\\
\end{tabular}

As the last case, the multiplication would still have to be enumerated, for which every arbitrary point is its own image.

The kind of the substitution is preserved under every transformation.

If a point remains unchanged by a substitution $A$, then it also remains unchanged under every repetition of $A$, hence under $A^2, A^3,\dots$. The poles of $A$ are therefore always found among the points that are their own images under the mapping by $A^2$. But the case may well occur that, for example, $A$ has three poles, while a power, $A^k$, possesses a principal axis. Then two of the poles of $A$ must lie on the principal axis of $A^k$, and the third pole of $A$ is the single pole of $A^k$. For if $A$ is a substitution with three poles, then, by placing the poles at the vertices of the coordinate triangle, we can put it in the normal form
\[
A=\begin{pmatrix}\alpha, & 0, & 0\\ 0, & \beta, & 0\\ 0, & 0, & \gamma\end{pmatrix}
\]
in which $\alpha,\beta,\gamma$ are different from one another. If now $\beta^k=\gamma^k$, so that $\beta$ differs from $\gamma$ by a $k$th root of unity as a factor, then $A^k$ has the line $x_1=0$ for its axis, and the opposite vertex of the coordinate triangle for its pole.

If $A$ is of finite degree $\mu$, then $\alpha,\beta,\gamma$ are $\mu$th roots of unity. If $A^k$ is the lowest power of $A$ that has a principal axis, then $\beta:\gamma$ must be a primitive $k$th and at the same time $\mu$th root of unity, and consequently $k$ must be a proper divisor of $\mu$; if $k$ is relatively prime to $\mu$, then $A^k$ has the same three poles as $A$.

\begin{center}
\S.\,117.

\gesperrt{Application to the group $G_{168}$. Sevenfold poles.}
\end{center}

We shall now seek out the poles and axes of the substitutions of our group $G_{168}$.
